\documentclass[10pt]{extarticle}
\usepackage[a4paper,margin=1.5cm]{geometry}
\usepackage[T1]{fontenc}
\usepackage{amsmath,amssymb}
\usepackage{graphicx}
\usepackage[export]{adjustbox}
\usepackage{multicol}
\usepackage{enumitem}
\usepackage{booktabs,array,tabularx}
\setlength{\columnsep}{0.7cm}
\setlength{\parindent}{0pt}
\setlength{\parskip}{3pt}
\pagestyle{empty}
\makeatletter
% figure registry (figdefs.tex is generated by the extract step)
\newcommand{\qfigdef}[4]{\@namedef{qf@file@#1}{#2}\@namedef{qf@w@#1}{#3}\@namedef{qf@wide@#1}{#4}}
\newcommand{\qf@check}[1]{\@ifundefined{qf@file@#1}{\PackageError{score2}{Unknown figure id #1}{}}{}}
\newcommand{\qf@img}[4][]{\includegraphics[width=#2\dimexpr\@nameuse{qf@w@#4}pt\relax,max width=#3,max height=0.45\textheight,#1]{\@nameuse{qf@file@#4}}}
% Full-width questions: finalize wraps every question that contains a wide figure in \qfullbegin...\qfullend,
% so the WHOLE question (stem, figures, options) is set across the page in reading order, then the two columns
% resume. The question's head - its number, stem and first figure (\qfullsplit marks the end of the head) -
% is kept on one page, so a number is never left at a page foot with its figure overleaf; the rest may break
% across pages normally, so no big gaps are left. Without a \qfullsplit, the whole question is kept together
% if it is under 45% of a page. Never write these yourself.
\newif\ifqf@full
\newif\ifqf@boxed
\newsavebox{\qf@box}
\newcommand{\qf@open}{\global\qf@fulltrue\par\global\qf@boxedtrue
  \setbox\qf@box\vbox\bgroup\hsize\textwidth\linewidth\textwidth}
\newcommand{\qf@emit}[1]{\par\egroup\global\qf@boxedfalse
  \ifdim\dimexpr\ht\qf@box+\dp\qf@box\relax<#1\textheight\noindent\usebox{\qf@box}\par\else\unvbox\qf@box\fi}
\newcommand{\qfullsplit}{\ifqf@boxed\qf@emit{0.85}\fi}
\newcommand{\qf@close}{\ifqf@boxed\qf@emit{0.45}\fi\global\qf@fullfalse}
\newcommand{\qfullbegin}{\end{multicols}\qf@open}
\newcommand{\qfullend}{\qf@close\begin{multicols}{2}}
\newcommand{\qfullstart}{\qf@open}
\newcommand{\qfulldone}{\qf@close}
% \qfig[scale]{id}: block figure, centred, directly where it is written (size set by the tool).
% Fallback if a wide figure is not inside a full-width question: break out of the columns in place - no floats.
\newcommand{\qfig}[2][1]{\qf@check{#2}%
  \ifqf@full\qf@blk{#1}{#2}\else
  \ifnum\@nameuse{qf@wide@#2}=1 \expandafter\@firstoftwo\else\expandafter\@secondoftwo\fi
  {\end{multicols}\begin{center}\qf@img{#1}{\textwidth}{#2}\end{center}\begin{multicols}{2}}%
  {\qf@blk{#1}{#2}}\fi}
\newcommand{\qf@blk}[2]{\par\smallskip{\centering\qf@img{#1}{\linewidth}{#2}\par}\smallskip}
% \qfigi[scale]{id}: inline figure (inside option lists / tables / minipages); never breaks columns.
\newcommand{\qfigi}[2][1]{\qf@check{#2}\qf@img[valign=c]{#1}{\linewidth}{#2}}
\makeatother
% structure macros - lint reads these to check completeness and figure placement
\newcommand{\papersec}[1]{\par\bigskip{\centering\Large\bfseries PAPER-#1\par}\medskip}
\newcommand{\partsec}[2]{\par\medskip\noindent{\large\bfseries PART-#1: #2}\par\smallskip}
\newcommand{\sect}[1]{\par\smallskip\noindent\textbf{#1}\par}
\newcommand{\qn}[1]{\par\smallskip\noindent\textbf{#1.}~}
\newcommand{\paraq}[2]{\par\smallskip\noindent\textbf{Paragraph for Questions #1 and #2}\par}
\newenvironment{qoptions}{\begin{itemize}[leftmargin=1.8em,itemsep=1pt,topsep=2pt,label={}]}{\end{itemize}}
\newcommand{\opt}[2]{\item[(#1)] #2}
% figure registry (generated by score2_tools.py finalize)
\qfigdef{p1_phy_s1_q03_1}{Enthusiast Pre-Score - Test 02 (11 Nov 2020)-img/p1_phy_s1_q03_1.png}{288.0}{1}
\qfigdef{p1_phy_s2_q05_1}{Enthusiast Pre-Score - Test 02 (11 Nov 2020)-img/p1_phy_s2_q05_1.png}{97.1}{0}
\qfigdef{p1_phy_s2_q06_1}{Enthusiast Pre-Score - Test 02 (11 Nov 2020)-img/p1_phy_s2_q06_1.png}{162.5}{0}
\qfigdef{p1_phy_s2_q07_1}{Enthusiast Pre-Score - Test 02 (11 Nov 2020)-img/p1_phy_s2_q07_1.png}{208.8}{0}
\qfigdef{p1_phy_s2_q08_1}{Enthusiast Pre-Score - Test 02 (11 Nov 2020)-img/p1_phy_s2_q08_1.png}{138.7}{0}
\qfigdef{p1_phy_s2_q11_1}{Enthusiast Pre-Score - Test 02 (11 Nov 2020)-img/p1_phy_s2_q11_1.png}{208.2}{0}
\qfigdef{p1_phy_s2_q12_1}{Enthusiast Pre-Score - Test 02 (11 Nov 2020)-img/p1_phy_s2_q12_1.png}{96.6}{0}
\qfigdef{p1_phy_s3_q01_1}{Enthusiast Pre-Score - Test 02 (11 Nov 2020)-img/p1_phy_s3_q01_1.png}{27.7}{0}
\qfigdef{p1_phy_s3_q04_1}{Enthusiast Pre-Score - Test 02 (11 Nov 2020)-img/p1_phy_s3_q04_1.png}{184.2}{0}
\qfigdef{p1_phy_s3_q05_1}{Enthusiast Pre-Score - Test 02 (11 Nov 2020)-img/p1_phy_s3_q05_1.png}{118.5}{0}
\qfigdef{p1_phy_s3_q06_1}{Enthusiast Pre-Score - Test 02 (11 Nov 2020)-img/p1_phy_s3_q06_1.png}{126.4}{0}
\qfigdef{p1_chem_s1_q03_1}{Enthusiast Pre-Score - Test 02 (11 Nov 2020)-img/p1_chem_s1_q03_1.png}{97.7}{0}
\qfigdef{p1_chem_s2_q05_A}{Enthusiast Pre-Score - Test 02 (11 Nov 2020)-img/p1_chem_s2_q05_A.png}{67.0}{0}
\qfigdef{p1_chem_s2_q05_B}{Enthusiast Pre-Score - Test 02 (11 Nov 2020)-img/p1_chem_s2_q05_B.png}{75.9}{0}
\qfigdef{p1_chem_s2_q05_C}{Enthusiast Pre-Score - Test 02 (11 Nov 2020)-img/p1_chem_s2_q05_C.png}{57.6}{0}
\qfigdef{p1_chem_s2_q05_D}{Enthusiast Pre-Score - Test 02 (11 Nov 2020)-img/p1_chem_s2_q05_D.png}{61.5}{0}
\qfigdef{p1_chem_s2_q06_1}{Enthusiast Pre-Score - Test 02 (11 Nov 2020)-img/p1_chem_s2_q06_1.png}{379.8}{1}
\qfigdef{p1_chem_s2_q11_1}{Enthusiast Pre-Score - Test 02 (11 Nov 2020)-img/p1_chem_s2_q11_1.png}{239.8}{0}
\qfigdef{p1_chem_s2_q11_A}{Enthusiast Pre-Score - Test 02 (11 Nov 2020)-img/p1_chem_s2_q11_A.png}{87.9}{0}
\qfigdef{p1_chem_s2_q11_B}{Enthusiast Pre-Score - Test 02 (11 Nov 2020)-img/p1_chem_s2_q11_B.png}{77.3}{0}
\qfigdef{p1_chem_s2_q11_C}{Enthusiast Pre-Score - Test 02 (11 Nov 2020)-img/p1_chem_s2_q11_C.png}{69.8}{0}
\qfigdef{p1_chem_s2_q11_D}{Enthusiast Pre-Score - Test 02 (11 Nov 2020)-img/p1_chem_s2_q11_D.png}{52.4}{0}
\qfigdef{p1_chem_s2_q12_1}{Enthusiast Pre-Score - Test 02 (11 Nov 2020)-img/p1_chem_s2_q12_1.png}{108.2}{0}
\qfigdef{p1_chem_s2_q12_A}{Enthusiast Pre-Score - Test 02 (11 Nov 2020)-img/p1_chem_s2_q12_A.png}{153.2}{0}
\qfigdef{p1_chem_s2_q12_B}{Enthusiast Pre-Score - Test 02 (11 Nov 2020)-img/p1_chem_s2_q12_B.png}{120.5}{0}
\qfigdef{p1_chem_s2_q12_C}{Enthusiast Pre-Score - Test 02 (11 Nov 2020)-img/p1_chem_s2_q12_C.png}{194.4}{0}
\qfigdef{p1_chem_s2_q12_3}{Enthusiast Pre-Score - Test 02 (11 Nov 2020)-img/p1_chem_s2_q12_3.png}{147.3}{0}
\qfigdef{p1_chem_s3_q03_1}{Enthusiast Pre-Score - Test 02 (11 Nov 2020)-img/p1_chem_s3_q03_1.png}{274.0}{1}
\qfigdef{p1_chem_s3_q04_1}{Enthusiast Pre-Score - Test 02 (11 Nov 2020)-img/p1_chem_s3_q04_1.png}{272.9}{1}
\qfigdef{p2_phy_s1_q02_1}{Enthusiast Pre-Score - Test 02 (11 Nov 2020)-img/p2_phy_s1_q02_1.png}{84.2}{0}
\qfigdef{p2_phy_s1_q03_1}{Enthusiast Pre-Score - Test 02 (11 Nov 2020)-img/p2_phy_s1_q03_1.png}{93.6}{0}
\qfigdef{p2_phy_s1_q04_1}{Enthusiast Pre-Score - Test 02 (11 Nov 2020)-img/p2_phy_s1_q04_1.png}{223.5}{0}
\qfigdef{p2_phy_s1_q07_1}{Enthusiast Pre-Score - Test 02 (11 Nov 2020)-img/p2_phy_s1_q07_1.png}{137.0}{0}
\qfigdef{p2_phy_s1_q07_2}{Enthusiast Pre-Score - Test 02 (11 Nov 2020)-img/p2_phy_s1_q07_2.png}{73.5}{0}
\qfigdef{p2_phy_s1_q07_3}{Enthusiast Pre-Score - Test 02 (11 Nov 2020)-img/p2_phy_s1_q07_3.png}{73.5}{0}
\qfigdef{p2_phy_s1_q08_1}{Enthusiast Pre-Score - Test 02 (11 Nov 2020)-img/p2_phy_s1_q08_1.png}{115.0}{0}
\qfigdef{p2_phy_s2_q09_1}{Enthusiast Pre-Score - Test 02 (11 Nov 2020)-img/p2_phy_s2_q09_1.png}{107.4}{0}
\qfigdef{p2_phy_s2_q10_1}{Enthusiast Pre-Score - Test 02 (11 Nov 2020)-img/p2_phy_s2_q10_1.png}{133.3}{0}
\qfigdef{p2_phy_s2_q11_1}{Enthusiast Pre-Score - Test 02 (11 Nov 2020)-img/p2_phy_s2_q11_1.png}{94.7}{0}
\qfigdef{p2_phy_s2_q12_1}{Enthusiast Pre-Score - Test 02 (11 Nov 2020)-img/p2_phy_s2_q12_1.png}{107.1}{0}
\qfigdef{p2_phy_s3_q02_1}{Enthusiast Pre-Score - Test 02 (11 Nov 2020)-img/p2_phy_s3_q02_1.png}{89.6}{0}
\qfigdef{p2_phy_s3_q03_1}{Enthusiast Pre-Score - Test 02 (11 Nov 2020)-img/p2_phy_s3_q03_1.png}{151.6}{0}
\qfigdef{p2_phy_s3_q06_1}{Enthusiast Pre-Score - Test 02 (11 Nov 2020)-img/p2_phy_s3_q06_1.png}{165.2}{0}
\qfigdef{p2_chem_s1_q03_1}{Enthusiast Pre-Score - Test 02 (11 Nov 2020)-img/p2_chem_s1_q03_1.png}{177.4}{0}
\qfigdef{p2_chem_s1_q03_A}{Enthusiast Pre-Score - Test 02 (11 Nov 2020)-img/p2_chem_s1_q03_A.png}{53.9}{0}
\qfigdef{p2_chem_s1_q03_B}{Enthusiast Pre-Score - Test 02 (11 Nov 2020)-img/p2_chem_s1_q03_B.png}{53.9}{0}
\qfigdef{p2_chem_s1_q03_C}{Enthusiast Pre-Score - Test 02 (11 Nov 2020)-img/p2_chem_s1_q03_C.png}{84.9}{0}
\qfigdef{p2_chem_s1_q03_D}{Enthusiast Pre-Score - Test 02 (11 Nov 2020)-img/p2_chem_s1_q03_D.png}{84.2}{0}
\qfigdef{p2_chem_s1_q04_1}{Enthusiast Pre-Score - Test 02 (11 Nov 2020)-img/p2_chem_s1_q04_1.png}{276.2}{1}
\qfigdef{p2_chem_s1_q06_1}{Enthusiast Pre-Score - Test 02 (11 Nov 2020)-img/p2_chem_s1_q06_1.png}{222.4}{0}
\qfigdef{p2_chem_s3_q03_1}{Enthusiast Pre-Score - Test 02 (11 Nov 2020)-img/p2_chem_s3_q03_1.png}{209.6}{0}
\raggedcolumns
\makeatletter\@beginparpenalty=10000\makeatother
\newcolumntype{L}[1]{>{\raggedright\arraybackslash}p{#1}}
\begin{document}
\begin{center}{\Large\bfseries Enthusiast Pre-Score --- Test 02 (11 Nov 2020)}\end{center}
\begin{multicols}{2}
\raggedcolumns
\papersec{1}
\partsec{1}{PHYSICS}
\sect{SECTION-I(i)}
\qn{1} An uncharged parallel plate capacitor having a dielectric of constant K is connected to a similar air cored parallel capacitor charged to a potential V. The two capacitors share charges and the common potential is V$'$. The dielectric constant K is :
\begin{qoptions}
\opt{A}{$\dfrac{V'-V}{V'+V}$}
\opt{B}{$\dfrac{V'-V}{V'}$}
\opt{C}{$\dfrac{V'-V}{V}$}
\opt{D}{$\dfrac{V-V'}{V'}$}
\end{qoptions}
\qn{2} In a meter bridge experiment the resistance of resistance box is $16\Omega$, which is inserted in right gap. The null point is obtained at 36 cm from the left end. The least count of meter scale is 1mm. The value of unknown resistance is
\begin{qoptions}
\opt{A}{$9\pm\dfrac{5}{128}\Omega$}
\opt{B}{$9\pm\dfrac{5}{250}\Omega$}
\opt{C}{$9\pm\dfrac{5}{512}\Omega$}
\opt{D}{$9\pm\dfrac{1}{2560}\Omega$}
\end{qoptions}
\qfullbegin
\qn{3} Consider the ladder network shown in figure. What should be the value of resistance R, so that effective resistance between A \& B becomes independent of number of elements in the combination?\par
\qfig{p1_phy_s1_q03_1}
\qfullsplit
\begin{qoptions}
\opt{A}{$2\,\Omega$}
\opt{B}{$4\,\Omega$}
\opt{C}{$8\,\Omega$}
\opt{D}{$16\,\Omega$}
\end{qoptions}
\qfullend
\qn{4} The diameter of a wire of length 50 cm is measured with the help of a screw gauge. The main scale reading is 1mm and circular scale reading is 25. Pitch of the screw gauge is 1mm and the total number of divisions on the circular scale is 100. This wire is used in an experiment of determination for the Young's modulus of material of a wire by Searle's method. The wire is placed under a tension of 50 N and elongation of wire is $\ell$ = 0.125 cm. The least count for measuring the length of wire is 0.01 cm and of elongation of the wire is 0.001 cm.\par
(A) The least count in the measurement of diameter of wire is 0.03 mm.\par
(B) The diameter of the wire is $(1.25\pm0.01)$ cm.\par
(C) The percentage error in calculation of Young's modulus is 2.4\%.\par
(D) The Young's modulus of wire is $(5.0\pm0.2)\times10^{10}$ N/m$^2$.
\sect{SECTION-I(ii)}
\qn{5} Shown below is an electric dipole with equal charges $+$Q and $-$Q separated by a distance d. The dipole is free to rotate or move. Consider the following information:\par
\qfig{p1_phy_s2_q05_1}
[I] The dipole does not feel a net torque about the axis of the dipole.\par
[II] The dipole feels a net force in the +x--direction.\par
[III] When rotated from its original orientation and released, the dipole moves back towards the original orientation.\par
Which of the following electric field, can justify the above information, if E$_0>0$.
\begin{qoptions}
\opt{A}{$\dfrac{E_0x}{d}\left(-\hat i\right)$}
\opt{B}{$\dfrac{E_0\left(x-\dfrac{3d}{4}\right)}{d}\left(-\hat i\right)$}
\opt{C}{$\dfrac{E_0\left(x-\dfrac{d}{4}\right)}{d}\left(-\hat i\right)$}
\opt{D}{$\dfrac{E_0x}{d}\left(\hat i\right)$}
\end{qoptions}
\qn{6} Two bulbs (B$_1$ : 25W, 100 V) and (B$_2$ : 100 W, 200 V) are connected in the circuit as shown in diagram : Choose \textbf{INCORRECT} option(s) :\par
\qfig{p1_phy_s2_q06_1}
\begin{qoptions}
\opt{A}{Heat lost per second in the circuit will be approximately 35 W.}
\opt{B}{Ratio of heat produced per second in both bulbs will be 1 : 1.}
\opt{C}{Ratio of heat produced in branch AB to branch CD will be 1 : 2.}
\opt{D}{Current drawn from the battery is 0.4 A.}
\end{qoptions}
\qn{7} In the given figure if A is connected with C and B is connected with D, then\par
\qfig{p1_phy_s2_q07_1}
\begin{qoptions}
\opt{A}{Charge flows in the circuit is 12 $\mu$C.}
\opt{B}{Heat produced in the circuit is 60 $\mu$J.}
\opt{C}{Charge flows in the circuit is 14 $\mu$C.}
\opt{D}{Heat produced in the circuit is 490 $\mu$J.}
\end{qoptions}
\qn{8} A tunnel is dug across the earth along a chord AB as shown. A particle is projected inside the tunnel from end `A' with an initial velocity $v_0=\dfrac{\sqrt{3g_0R}}{2}$, there is no friction between the tunnel and the particle. Then choose the correct option(s). (where g$_0$ = acceleration due to gravity on the surface of Earth, R = radius of earth).\par
\qfig{p1_phy_s2_q08_1}
\begin{qoptions}
\opt{A}{The time taken by the particle to move from end `A' to other end `B' of the tunnel is $\dfrac{\pi}{3}\sqrt{\dfrac{R}{g_0}}$.}
\opt{B}{The time taken by the particle to move from end `A' to other `B' of the tunnel is $\dfrac{2\pi}{3}\sqrt{\dfrac{R}{g_0}}$.}
\opt{C}{The maximum speed of the particle inside the tunnel is $\sqrt{g_0R}$}
\opt{D}{The maximum speed of the particle inside the tunnel is $\sqrt{2g_0R}$}
\end{qoptions}
\qn{9} A meter scale of mass 0.50 g is placed on smooth horizontal surface. The least count of scale is 0.1 cm. An earthworm of mass 0.25 g is on scale such that length of earthworm is between mark of 4.4 cm and 16.1 cm. The earthworm crawl along the length of meter scale such that its length changes during crawling (but centre of mass remains at geometrical centre of its body). After (10.0 $\pm$ 0.1) sec its length is between mark 20.0 cm and 36.5 cm. Then choose the \textbf{CORRECT} statement (s) :-
\begin{qoptions}
\opt{A}{With respect to ground magnitude of displacement of meter scale is 6.0 cm}
\opt{B}{Change in length of the earthworm is 4.8 cm}
\opt{C}{With respect to meter scale maximum error in measurement of change in length of the earthworm is 0.2 cm}
\opt{D}{With respect to meter scale maximum error in measurement of change in length of the earthworm is 0.4 cm}
\end{qoptions}
\qn{10} An earth satellite is revolving in a circular orbit of radius a with a velocity v$_0$. Mass of earth is M. A gun in the satellite is directly aimed toward earth. A bullet is fired from the gun with muzzle velocity $\dfrac{v_0}{2}$. Assume that mass of the satellite is very-very large with respect to the mass of the bullet. Choose the \textbf{\emph{CORRECT}} option(s) :
\begin{qoptions}
\opt{A}{Maximum distance of bullet from the centre of earth is $\dfrac{4a}{3}$.}
\opt{B}{Minimum speed of bullet (relative to earth) is $\dfrac{v_0}{2}$.}
\opt{C}{Minimum speed of bullet (relative to earth) is $\dfrac{3v_0}{4}$.}
\opt{D}{Time period of motion of bullet around earth is $T=\dfrac{2\pi}{\sqrt{GM}}\left(\dfrac{4a}{3}\right)^{3/2}$.}
\end{qoptions}
\qn{11} In the given circuit the potential difference between a \& b is 300 V. The resistors are of 100 $\Omega$ each.\par
\qfig{p1_phy_s2_q11_1}
\begin{qoptions}
\opt{A}{If the upper meter M$_1$ is ideal voltmeter \& lower is ideal ammeter M$_2$, reading of voltmeter is 300 V.}
\opt{B}{If both meters are ideal voltmeters, they show 200 V each.}
\opt{C}{If both meters are ideal ammeters, they read 6A each.}
\opt{D}{If the upper meter M$_1$ is ideal voltmeter and lower is ideal ammeter M$_2$, reading of ammeter is 3A.}
\end{qoptions}
\qn{12} A metallic sphere of radius 1m has a spherical cavity of radius $\dfrac12$m in it as shown. The centers are at a distance of 20 cm from each other. A charge of q = + 5$\mu$C is kept in the cavity on the line joining the centers at a distance of 40 cm from centre of the original sphere. If net charge on the metallic sphere is +4$\mu$C, then\par
\qfig{p1_phy_s2_q12_1}
\begin{qoptions}
\opt{A}{Potential at centre A is 103.5 kV.}
\opt{B}{Potential at centre B of cavity is 216 kV.}
\opt{C}{The electric field at B is 1125 kN/C.}
\opt{D}{The electric field outside the sphere does not change if the position of +q inside cavity changes.}
\end{qoptions}
\sect{SECTION-II}
\qn{1} In the figure shown S is a large non-conducting sheet of uniform charge density $\sigma$. A rod R of length $\ell$ and mass `m' is parallel to the sheet and hinged at its mid-point. The linear charge densities on the upper and lower half of the rod are shown in the figure. Find the angular acceleration (in radian/s$^2$) of the rod just after it is released.\par
(Take : $\lambda=\dfrac1\pi\,\mu$C/m, $\sigma$ = 1mC/m$^2$, m = 2kg)\par
\qfig{p1_phy_s3_q01_1}
\qn{2} A particle of mass (m) is revolving around a planet of mass `M' in circular orbit of radius `3R'. It is given $\dfrac{GMm}{10aR}$ energy, so that now it moves in elliptical orbit of semi major axis `5R'. Find `a'.
\qn{3} The pitch of a screw gauge is 1mm and its cap is divided into 100 divisions. When nothing is placed between its studs, the zero of circular scale is 4 divisions above the reference line. When a wire is placed between its studs, the main scale reading is 2mm and 46$^{\text{th}}$ division of circular scale coincides with the reference line. The length of wire is 4.76 cm. Find the curved surface area (in cm$^2$) of wire in true significant figures.
\qn{4} In the circuit shown, the switch `S' is closed at t = 0. Then find the energy (in milli Joule) stored in 5 $\mu$F capacitor in steady state.\par
\qfig{p1_phy_s3_q04_1}
\qn{5} Earth revolves around Sun with time period of `T = 365 days'. In the figure shown below, find the time taken (in days) by the Earth to move from A to B in the path ACB if the eccentricity of ellipse is 0.5.\par
\qfig{p1_phy_s3_q05_1}
\qn{6} In the circuit shown in figure, initially the terminals A and P are connected together to charge 4$\mu$F capacitor. Then battery is disconnected, and the terminals P and B are connected together. Then the ratio of final charge and initial charge on 4 $\mu$F capacitor is :\par
\qfig{p1_phy_s3_q06_1}
\partsec{2}{CHEMISTRY}
\sect{SECTION-I(i)}
\qn{1} (1) Bottle-1 (Colourless solution) + Pb(OAc)$_2$ solution $\longrightarrow$ white ppt. which dissolves only on strong heating ($\sim$85-90$^\circ$C)\par
(2) Bottle-2 (Coloured solution) + Pb(OAc)$_2$ solution $\longrightarrow$ yellow ppt. insoluble in acidic acid.\par
(3) Solid salts in bottle-1 and bottle-2 are mixed together and taking that mixture, flame test is performed. Voluminous yellow flame is observed with naked eye and Crimson red flame is observed through double blue glass.\par
Which of the following statement is \textbf{INCORRECT}?
\begin{qoptions}
\opt{A}{The anion present in bottle-2 is $\mathrm{CrO_4^{2-}}$}
\opt{B}{The cation present in bottle-1 may be Na$^+$}
\opt{C}{The cation present in bottle-2 may be Na$^+$}
\opt{D}{The anion present in bottle-1 is Cl$^-$}
\end{qoptions}
\qn{2} In the group analysis of basic radicals, the colour of the ppt. obtained in group-III test is reddish brown. The oxidation state of cation present in the original salt may be.
\begin{qoptions}
\opt{A}{+2}
\opt{B}{+3}
\opt{C}{+2 and +3}
\opt{D}{All of these}
\end{qoptions}
\qn{3} The given graph is for salt hydrolysis (monovalent type) between pH of solution vs. log c (where c = concentration of salt)\par
\qfig{p1_chem_s1_q03_1}
If 45 ml, 0.2 M of above salt is titrated with 60 ml, 0.1 M of NaOH, then pH of resulting solution will be :
\begin{qoptions}
\opt{A}{7.7}
\opt{B}{8.3}
\opt{C}{5.7}
\opt{D}{6.3}
\end{qoptions}
\qn{4} 4.16 kg of an equimolar mixture of ferrous oxalate (Molar mass = 144 g mole$^{-1}$) and ferric oxalate (Molar mass = 376 g mole$^{-1}$) require `X' moles of K$_2$Cr$_2$O$_7$ in acidic medium, then value of `X' ?
\begin{qoptions}
\opt{A}{12}
\opt{B}{6}
\opt{C}{$\dfrac{16}{3}$}
\opt{D}{$\dfrac{32}{3}$}
\end{qoptions}
\sect{SECTION-I(ii)}
\qn{5} What is the most likely major final product when 1-methylcyclopentene is treated with the following sequence of reatents :\par
(1) BH$_3$ \quad (2) NaOH / H$_2$O$_2$ \quad (3) PCC / CH$_2$Cl$_2$ \quad (4) Hydrazine \quad (5) KOH / heat
\begin{qoptions}
\opt{A}{\qfigi{p1_chem_s2_q05_A}}
\opt{B}{\qfigi{p1_chem_s2_q05_B}}
\opt{C}{\qfigi{p1_chem_s2_q05_C}}
\opt{D}{\qfigi{p1_chem_s2_q05_D}}
\end{qoptions}
\qfullbegin
\qn{6} Which of the following are present in the equilibrium mixture of acetone and methanol at pH 3 ?\par
\qfig{p1_chem_s2_q06_1}
\qfullsplit
\begin{qoptions}
\opt{A}{i, ii, iii}
\opt{B}{i, iii, iv, v}
\opt{C}{i, iv}
\opt{D}{ii, iii, iv}
\end{qoptions}
\qfullend
\qn{7} Which of the following statements is/are \textbf{CORRECT} for iodide of Millon's base?
\begin{qoptions}
\opt{A}{The formation reaction of it is a redox reaction}
\opt{B}{The number atoms per molecule is 7}
\opt{C}{The colour of it is brown}
\opt{D}{The test associated to it is used to detect NH$_3$ gas}
\end{qoptions}
\qn{8} Salt solution `A' is treated with H$_2$S in presence of dil. HCl and resulting a yellow ppt. which dissolves in YAS. The cation present in the salt may be.
\begin{qoptions}
\opt{A}{Sn$^{2+}$}
\opt{B}{Sn$^{4+}$}
\opt{C}{Cd$^{2+}$}
\opt{D}{As$^{+5}$}
\end{qoptions}
\qn{9} 51 gm sample of a double salt Fe$_x$ Al.H$_y$(C$_2$O$_4$)$_z$.nH$_2$O (molar mass = 510 gm) is dissolved in water to form 500 ml solution. 50 ml solution required 30 ml, 0.5 M Ba(OH)$_2$ and 180 ml, 0.1M KMnO$_4$ in acidic medium separately in two different experiments.\par
[At. mass : Fe = 56, Al = 27].\par
Select correct statement(s) [take all Fe to be present in (+ 2) oxidation state in the double salt]
\begin{qoptions}
\opt{A}{(x + y + z $-$ n) is equal 4}
\opt{B}{x = y}
\opt{C}{Number of oxalate ions per molecule is 3}
\opt{D}{Value of `n' is 4}
\end{qoptions}
\qn{10} Select CORRECT statement(s)
\begin{qoptions}
\opt{A}{S.I. unit of molar conductance is S-m$^2$/mol}
\opt{B}{Molar conductance of electrolytic solution decreases with increasing temperature.}
\opt{C}{Conductance of electrolyte in a solution depends on distance between electrodes}
\opt{D}{Conductivity increases with increase in concentration of electrolyte in solution.}
\end{qoptions}
\qn{11} Identify reactant in following sequence?\par
\qfig{p1_chem_s2_q11_1}
\begin{qoptions}
\opt{A}{\qfigi{p1_chem_s2_q11_A}}
\opt{B}{\qfigi{p1_chem_s2_q11_B}}
\opt{C}{\qfigi{p1_chem_s2_q11_C}}
\opt{D}{\qfigi{p1_chem_s2_q11_D}}
\end{qoptions}
\qn{12} Using \qfigi{p1_chem_s2_q12_1} as a reactant identify reactions correctly matched with their major products?\par
\qfig{p1_chem_s2_q12_A}
\qfig{p1_chem_s2_q12_B}
\qfig{p1_chem_s2_q12_C}
\qfig{p1_chem_s2_q12_3}
\sect{SECTION-II}
\qn{1} Find the ratio of number of incorrect to correct statements from the following.\par
(i) In the layer test, the organic layer is used to intensify the observation.\par
(ii) Layer test is the confimation test for Br$^-$ and Cl$^-$.\par
(iii) K$_4$[Fe(CN)$_6$] can be used to distinguish between Ca$^{2+}$ and Sr$^{+2}$\par
(iv) The addition of Na$_2$C$_2$O$_4$ solution followed by acetic acid can distinguish between Ca$^{2+}$ and Ba$^{2+}$\par
(v) K$^+$ and NH$_4^+$ can be distinguished by the addition of H$_2$[PtCl$_6$].
\qn{2} What is the sum of oxidation number and coordination number of central metal atom of complex formed as yellow precipitate when CoCl$_2$ reacts with excess of potassium nitrite ?
\qfullbegin
\qn{3} Identify compounds give positive test with NaHSO$_3$?\par
\qfig{p1_chem_s3_q03_1}
\qn{4} Identify compounds give white ppt. with Ph--SO$_2$Cl which is soluble in aqueous NaOH ?\par
\qfig{p1_chem_s3_q04_1}
\qfullend
\qn{5} How many seconds would it take for the pH of solution to rise to a value of 10, if 60 ml of NaI solution is electrolysed with a current of 0.1 A
\qn{6} If a solution contains 0.2 M each of Mg$^{+2}$ and Mn$^{+2}$ ions, to what pH must this solution be made so as to maximise precipitation of Mn(OH)$_2$ and almost no precipitation of Mg(OH)$_2$\par
Ksp [Mn(OH)$_2$] = 1.6 $\times$ 10$^{-13}$ ; Ksp [Mg(OH)$_2$] = 1.2 $\times$ 10$^{-11}$\par
Given (log 6 = 0.78 and log 8 = 0.9)
\partsec{3}{MATHEMATICS}
\sect{SECTION-I(i)}
\qn{1} In a triangle ABC, with usual notation $\angle$ABC = 45$^\circ$, Point D is on BC so that 2BD = CD and $\angle$DAB = 75$^\circ$ then $\angle$ACB equals to
\begin{qoptions}
\opt{A}{15$^\circ$}
\opt{B}{$\angle$A $-$ $\angle$B}
\opt{C}{30$^\circ$}
\opt{D}{$\dfrac{\angle A-\angle B}{2}$}
\end{qoptions}
\qn{2} In $\Delta$ABC if B = 2A then $\dfrac{b^2c+a^3}{b^2+c^2}$ is equal to (symbols have their usual meaning)
\begin{qoptions}
\opt{A}{a}
\opt{B}{b}
\opt{C}{c}
\opt{D}{a + b $-$ c}
\end{qoptions}
\qn{3} The solution of $4\sin^2x+\tan^2x+\operatorname{cosec}^2x+\cot^2x-6=0$, is (where n $\in$ I) -
\begin{qoptions}
\opt{A}{$n\pi\pm\dfrac{\pi}{4}$}
\opt{B}{$2n\pi\pm\dfrac{\pi}{4}$}
\opt{C}{$n\pi+\dfrac{\pi}{3}$}
\opt{D}{$n\pi-\dfrac{\pi}{6}$}
\end{qoptions}
\qn{4} If the inequation $\sqrt{1-x^2}\ge a-|x|$ holds good for all x $\in$ [$-$1, 1], then the value of `a' cannot be
\begin{qoptions}
\opt{A}{$\dfrac12$}
\opt{B}{1}
\opt{C}{$\dfrac43$}
\opt{D}{$\dfrac34$}
\end{qoptions}
\sect{SECTION-I(ii)}
\qn{5} Let $\cos20^\circ+2\cos40^\circ=a$\par
$\sin20^\circ+2\sin40^\circ=b$\par
then identify the correct statement(s)
\begin{qoptions}
\opt{A}{$a^2+b^2>\dfrac{42}{5}$}
\opt{B}{$a^2+b^2<9$}
\opt{C}{$\dfrac{\sin25^\circ+2\sin5^\circ}{\cos25^\circ+2\cos5^\circ}=\dfrac{a-b}{a+b}$}
\opt{D}{$\dfrac{\sin25^\circ+2\sin5^\circ}{\cos25^\circ+2\cos5^\circ}=\dfrac{2(a-b)}{a+b}$}
\end{qoptions}
\qn{6} Given x and y are real numbers satisfying $2\sin x\sin y+3\cos y+6\cos x\sin y=7$, then which of the following can be true-
\begin{qoptions}
\opt{A}{$\dfrac{\cos x+\sin x}{\cos x-\sin x}=2$}
\opt{B}{$2\tan^2x+3\sec^2x=\dfrac{79}{9}$}
\opt{C}{$\tan^2x+\sec^2y=\dfrac{50}{9}$}
\opt{D}{$\tan^2x+\sec^2y=\dfrac{49}{9}$}
\end{qoptions}
\qn{7} Let $S_n(\theta)=\displaystyle\sum_{r=0}^{n}\dfrac{2\sec\left(2^{r+2}\theta\right)}{\cot\left(2^r\theta\right)-\tan\left(2^r\theta\right)}$. Identify the correct statement(s) -
\begin{qoptions}
\opt{A}{$S_3\left(\dfrac\pi6\right)=0$}
\opt{B}{$S_5\left(\dfrac\pi3\right)=-\sqrt3$}
\opt{C}{$S_5\left(\dfrac\pi3\right)=0$}
\opt{D}{$S_3\left(\dfrac\pi6\right)=-2\sqrt3$}
\end{qoptions}
\qn{8} Let $\sqrt{2\cos^2\theta-4\cos\theta+6}+\sqrt{19-6\cos\theta+3\cos^2\theta}=4+\sin^2\theta+2\cos\theta$ where $0\le\theta\le4\pi$ then :
\begin{qoptions}
\opt{A}{Number of solutions of the equation is 2}
\opt{B}{Number of solution the equation is 3}
\opt{C}{Sum of all solutions is $3\pi$}
\opt{D}{Sum of all solutions is $6\pi$}
\end{qoptions}
\qn{9} In a triangle ABC length of medians AD and CE are 18 and 27 respectively. CE is extended to intersect its circumcircle of ABC at F. If AB = 24 and G is centroid of the triangle, then-
\begin{qoptions}
\opt{A}{b = $3\sqrt{70}$}
\opt{B}{a = $6\sqrt{31}$}
\opt{C}{$\Delta$AGE = $\dfrac{27}{4}\sqrt{55}$}
\opt{D}{EF = $\dfrac{16}{3}$}
\end{qoptions}
\qn{10} In an isosceles $\Delta$ABC with integral sides, AB = 7 cm and $\cos A=\dfrac27$, then which of the following is/are true
\begin{qoptions}
\opt{A}{Area of $\Delta$ABC is $6\sqrt5$ cm$^2$}
\opt{B}{Area of $\Delta$ABC is $\dfrac{1029}{32}\sqrt5$ cm$^2$}
\opt{C}{Circumradius of $\Delta$ABC is $\dfrac{49\sqrt5}{30}$ cm}
\opt{D}{Circumradius of $\Delta$ABC is $\dfrac{49\sqrt5}{12}$ cm}
\end{qoptions}
\qn{11} In $\Delta$ABC, if a,b,c are in A.P., then which of the following are always in A.P. (symbols have their usual meaning)
\begin{qoptions}
\opt{A}{sinA, sinB, sinC}
\opt{B}{$\dfrac1{r_1},\dfrac1{r_2},\dfrac1{r_3}$}
\opt{C}{$\cot\dfrac A2,\cot\dfrac B2,\cot\dfrac C2$}
\opt{D}{Distance of circumcentre from sides BC, CA and AB respectively.}
\end{qoptions}
\qn{12} Which of the following statement(s) is(are) correct ?
\begin{qoptions}
\opt{A}{Number of solutions of the equation cos(tanx) = sin(tanx) in x $\in\left(0,\dfrac\pi2\right)$ is infinite.}
\opt{B}{If the equation in x, sinx + cos(k + x) + cos(k $-$ x) = 2 has a real solution, then $|\sin k|\le\dfrac12$.}
\opt{C}{The equation 4sin2x + cosx = 5 has no real solution.}
\opt{D}{Number of solutions of the equation sinx.cos4x.sin5x = $\dfrac{-1}{2}$ in x $\in\left(-\pi,\dfrac{3\pi}2\right]$ is two.}
\end{qoptions}
\sect{SECTION-II}
\qn{1} The value of $\tan3^\circ\left(\tan3^\circ+2\tan6^\circ+\right.$ $\left.4\tan12^\circ+8\cot24^\circ\right)$ is equal to
\qn{2} The value of $\sqrt3\tan50^\circ-4\sin50^\circ$ is-
\qn{3} If A,B $>$ 0 and $A+B=\dfrac\pi3$, then maximum value of tanA tanB is
\qn{4} In a $\Delta$ABC, if the internal angle bisector of angle A meeting BC at L and median BE are perpendicular, then $\dfrac{BE\cot\frac A2}{AL}$ is equal to
\qn{5} A cyclic quadrilateral ABCD of area $\dfrac{3\sqrt3}{4}$ is inscribed in a unit circle. If AB = 1, BD = $\sqrt3$, then perimeter of cyclic quadrilateral is
\qn{6} If $-x^2+3(\sin\theta)x-\cos\phi(x-\sin\theta)\le0\ \forall\ x\in R$, then number of ordered pair ($\theta$,$\phi$), $\theta,\phi\in[0,2\pi]$ is
\papersec{2}
\partsec{1}{PHYSICS}
\sect{SECTION-I(i)}
\qn{1} Two Earth's satellites move in a common plane along circular obrits. The orbital radius of one satellite is r while that of the other satellite is r $-$ $\Delta$r (Here $\Delta$r $\ll$ r).
\begin{qoptions}
\opt{A}{Time interval separates the periodic approaches of the satellites to each other over the minimum distance is $\dfrac{4\pi r^{5/2}}{3(GM)^{1/2}\Delta r}$.}
\opt{B}{Time interval separates the periodic approaches of the satellites to each other over the minimum distance is $\dfrac{2\pi r^{5/2}}{3(GM)^{1/2}\Delta r}$.}
\opt{C}{Angular velocity of approach between two satellites is $\dfrac{3(GM)^{1/2}\Delta r}{r^{5/2}}$.}
\opt{D}{Angular velocity of approach between two satellites is $\dfrac{3(GM)^{1/2}\Delta r}{2\,r^{5/2}}$.}
\end{qoptions}
\qn{2} A regular tetrahedron frame is made up of homogeneous resitance wire of uniform cross section. Current I enters vertex A through a long, straight wire directed towards the centre O of the tetrahedron and is conducted away through vertex B in the same way, as illustrated in the figure. Choose \textbf{CORRECT} statement(s) :-\par
\qfig{p2_phy_s1_q02_1}
\begin{qoptions}
\opt{A}{There is no current through wire DC because point D and C are equipotential.}
\opt{B}{Current through wire AB is I/2}
\opt{C}{Current through wire AC is I/4}
\opt{D}{Current through wire AD is I/4}
\end{qoptions}
\qn{3} Three identical, parallel conducting plates A, B and C are placed as shown. Switches S$_1$ and S$_2$ are open, and can connect A and C to earth when closed. $+$Q charge is given to B.\par
\qfig{p2_phy_s1_q03_1}
\begin{qoptions}
\opt{A}{If S$_1$ is closed with S$_2$ open, a charge of amount Q will pass through S$_1$}
\opt{B}{If S$_2$ is closed with S$_1$ open, a charge of amount Q will pass through S$_2$}
\opt{C}{If S$_1$ and S$_2$ are closed together, a charge of amount Q/3 will pass through S$_1$, and a charge of amount 2Q/3 will pass through S$_2$.}
\opt{D}{All the above statements are incorrect}
\end{qoptions}
\qn{4} Which of the following options are \textbf{\emph{CORRECT}} (ammeters are ideal).\par
\qfig{p2_phy_s1_q04_1}
\begin{qoptions}
\opt{A}{Reading of A$_1$ is 1.16 A}
\opt{B}{Reading of voltmeter v is 77.14 V}
\opt{C}{Reading of ammeter A$_2$ is 0.77 A}
\opt{D}{Reading of voltmeter is 20 V}
\end{qoptions}
\qn{5} Using screw gauge the diameter of the wire is found to be 5.00 mm. The length of wire is measured by using a scale and is found to be 50.0cm. If mass of wire is measured as 25g, then mark the \textbf{CORRECT} statement(s) (Take $\pi$ = 3.14).
\begin{qoptions}
\opt{A}{The density has to be computed upto 2 significant digits.}
\opt{B}{The least count of scale used to measure length of wire is 1 mm.}
\opt{C}{The density of wire is 2.5 g/cm$^3$}
\opt{D}{The least count of screw gauge is 0.01 mm}
\end{qoptions}
\qn{6} While measuring length of small iron rod, following five readings were obtained.\par
12.66 mm, 12.58 mm, 12.62 mm, 12.80 mm and 12.72 mm\par
The experimenter reported the result of the experiment based on the above readings.
\begin{qoptions}
\opt{A}{Number of significant figures in the length reported will be 3.}
\opt{B}{Number of significant figures in the length reported will be 4.}
\opt{C}{Number of significant figures in the error reported will be 1.}
\opt{D}{Number of significant figures in the error reported will be 2.}
\end{qoptions}
\qn{7} The diagram below shows some lines of equipotential in the region of an electric field.\par
\qfig{p2_phy_s1_q07_1}
\begin{qoptions}
\opt{A}{Graph which shows the magnitude $E$ of the electric field strength along the line XY is\par\qfig{p2_phy_s1_q07_2}}
\opt{B}{Graph which shows the magnitude $E$ of the electric field strength along the line XY is\par\qfig{p2_phy_s1_q07_3}}
\opt{C}{If an electron is released from rest (from point X) then its velocity increases \& its acceleration first remains constant and then decreases}
\opt{D}{If an electron is released from rest (from point X) then its velocity increases \& its acceleration first remains constant and then increases}
\end{qoptions}
\qn{8} In the figure below a neutral insulating disk of mass M and radius R is pivoted at the center. A small charge $+$q (having negligible mass) is fixed to one edge and a charge of $-$q (having negligible mass) is fixed to the opposite edge so that the line between the two charges makes a small angle $\theta$ with an applied, uniform electric field $\vec E$ as shown.\par
\qfig{p2_phy_s1_q08_1}
\begin{qoptions}
\opt{A}{The torque on the disk is q2RE sin$\theta$}
\opt{B}{The angular frequency with which the disk oscillates $\sqrt{\dfrac{4Eq}{MR}}$.}
\opt{C}{The angular frequency with which the disk oscillates $\sqrt{\dfrac{2Eq}{MR}}$.}
\opt{D}{The torque on the disk is RE sin$\theta$}
\end{qoptions}
\sect{SECTION-I(ii)}
\qn{9} \textbf{Answer the following by appropriately matching the lists based on the information given in the paragraph.}\par
Consider the situation shown in diagram. Both capacitors are having same capacitance C. Capacitors are initially uncharged. List-I represents value of r and List-II represents charge passing through r during complete process. Now match the followings :\par
\qfig{p2_phy_s2_q09_1}
Switch K is closed at time = 0. After a long time capacitors are fully charged.\par
\smallskip
\noindent{\footnotesize\setlength{\tabcolsep}{2pt}\begin{tabular}{|L{.95cm}|L{2.7cm}|L{.95cm}|L{3.1cm}|}\hline
\multicolumn{2}{|c|}{List-I} & \multicolumn{2}{c|}{List-II}\\\hline
(I) & r = 0 & (P) & $\dfrac{CV}{6}$\\\hline
(II) & r = R & (Q) & $\dfrac{CV}{5}$\\\hline
(III) & r = 2R & (R) & $\dfrac{CV}{4}$\\\hline
(IV) & r = 3R & (S) & $\dfrac{CV}{3}$\\\hline
 & & (T) & Zero\\\hline
 & & (U) & $\dfrac{CV}{2}$\\\hline
\end{tabular}}
\begin{qoptions}
\opt{A}{I $\to$ S, II $\to$ T, III $\to$ Q, IV $\to$ U}
\opt{B}{I $\to$ P, II $\to$ R, III $\to$ T, IV $\to$ S}
\opt{C}{I $\to$ P, II $\to$ T, III $\to$ Q, IV $\to$ T}
\opt{D}{I $\to$ S, II $\to$ R, III $\to$ Q, IV $\to$ P}
\end{qoptions}
\qn{10} \textbf{Answer the following by appropriately matching the lists based on the information given in the paragraph.}\par
Consider the situation shown in diagram. Both capacitors are having same capacitance C. Capacitors are initially uncharged. List-I represents value of r and List-II represents charge passing through r during complete process. Now match the followings :\par
\qfig{p2_phy_s2_q10_1}
Switch K is opened after reaching steady state.\par
\smallskip
\noindent{\footnotesize\setlength{\tabcolsep}{2pt}\begin{tabular}{|L{.95cm}|L{2.7cm}|L{.95cm}|L{3.1cm}|}\hline
\multicolumn{2}{|c|}{List-I} & \multicolumn{2}{c|}{List-II}\\\hline
(I) & r = 0 & (P) & $\dfrac{CV}{6}$\\\hline
(II) & r = R & (Q) & $\dfrac{CV}{5}$\\\hline
(III) & r = 2R & (R) & $\dfrac{CV}{4}$\\\hline
(IV) & r = 3R & (S) & $\dfrac{CV}{3}$\\\hline
 & & (T) & Zero\\\hline
 & & (U) & $\dfrac{CV}{2}$\\\hline
\end{tabular}}
\begin{qoptions}
\opt{A}{I $\to$ S, II $\to$ T, III $\to$ Q, IV $\to$ U}
\opt{B}{I $\to$ T, II $\to$ T, III $\to$ T, IV $\to$ T}
\opt{C}{I $\to$ P, II $\to$ P, III $\to$ R, IV $\to$ R}
\opt{D}{I $\to$ S, II $\to$ Q, III $\to$ R, IV $\to$ P}
\end{qoptions}
\qn{11} \textbf{Answer the following by appropriately matching the lists based on the information given in the paragraph.}\par
Match the entries on left side with entries on right side. Three identical metal plates with large surface areas are kept parallel to each other as shown in figure. The left most is given a charge Q, the right most a charge $-$2Q and the middle one remains neutral. Then\par
\qfig{p2_phy_s2_q11_1}
\smallskip
\noindent{\footnotesize\setlength{\tabcolsep}{2pt}\begin{tabular}{|L{.95cm}|L{2.7cm}|L{.95cm}|L{3.1cm}|}\hline
\multicolumn{2}{|c|}{List-I} & \multicolumn{2}{c|}{List-II}\\\hline
(I) & The charge appearing on outer surface of right most plate is & (P) & $+\dfrac Q2$\\\hline
(II) & The charge appearing on outer surface of left most plate is & (Q) & $-\dfrac Q2$\\\hline
(III) & The charge appearing on left surface of middle plate is & (R) & $-\dfrac{3Q}2$\\\hline
(IV) & The charge appearing on right surface of middle plate is & (S) & $\dfrac{3Q}2$\\\hline
 & & (T) & $+$Q\\\hline
 & & (U) & $-$ Q\\\hline
\end{tabular}}
\begin{qoptions}
\opt{A}{I $\to$ P, II $\to$ Q, III $\to$ R, IV $\to$ U}
\opt{B}{I $\to$ P, II $\to$ Q, III $\to$ R, IV $\to$ S}
\opt{C}{I $\to$ P, II $\to$ T, III $\to$ R, IV $\to$ U}
\opt{D}{I $\to$ Q, II $\to$ Q, III $\to$ R, IV $\to$ S}
\end{qoptions}
\qn{12} \textbf{Answer the following by appropriately matching the lists based on the information given in the paragraph.}\par
Match the entries on left side with entries on right side when left most plate is grounded. Then\par
\qfig{p2_phy_s2_q12_1}
\smallskip
\noindent{\footnotesize\setlength{\tabcolsep}{2pt}\begin{tabular}{|L{.95cm}|L{2.7cm}|L{.95cm}|L{3.1cm}|}\hline
\multicolumn{2}{|c|}{List-I} & \multicolumn{2}{c|}{List-II}\\\hline
(I) & The charge appearing on outer surface of right most plate is & (P) & 0\\\hline
(II) & The charge appearing on outer surface of left most plate is & (Q) & Q\\\hline
(III) & The charge appearing on left surface of middle plate is & (R) & $-$Q\\\hline
(IV) & The charge appearing on right surface of middle plate is & (S) & 2Q\\\hline
 & & (T) & $-$2Q\\\hline
 & & (U) & 3Q\\\hline
\end{tabular}}
\begin{qoptions}
\opt{A}{I $\to$ P, II $\to$ Q, III $\to$ R, IV $\to$ U}
\opt{B}{I $\to$ P, II $\to$ P, III $\to$ T, IV $\to$ U}
\opt{C}{I $\to$ P, II $\to$ P, III $\to$ T, IV $\to$ S}
\opt{D}{I $\to$ Q, II $\to$ Q, III $\to$ R, IV $\to$ S}
\end{qoptions}
\sect{SECTION-II}
\qn{1} A trinary star system which is a system of three stars orbiting around centre of mass of system has time period of 3 years, while the distance between any two stars of the system is 2 astronomical unit. All the three stars are identical and mass of the sun is M and total mass of this multiple star system is $\dfrac ab\times M$. Then find $\dfrac b3$? (a and b are lowest possible integers)\par
(1 astronomical unit is equal to distance between sun and earth, time period of rotation of earth around sun is 1year)
\qn{2} Figure shows two concentric spherical shall having radii 1 m and 2 m respectively. Outer spherical surface has charge $-3.2\times10^{-15}$ Coulomb. If inner surface is now earthed by closing the switch then x $\times$ 10$^{+y}$ number of electron move towards earth, then x + y will be (x, y $\in$ I, 0 $\le$ x,y $\le$ 9).\par
\qfig{p2_phy_s3_q02_1}
\qn{3} In an experiment with a potentiometer, V$_B$ = 10V, R is adjusted to be 50$\Omega$. A student wanting to measure voltage E$_1$ of a battery (approx. 8V) finds no null point possible. He then diminshes R to 10$\Omega$ and is able to locate the null point on the last (4$^{\text{th}}$) segment of the potentiometer. The resistance of the potentiometer wire ($R'$) will be $a<R'<b$. Find the value of $\left(\dfrac ba\right)$.\par
\qfig{p2_phy_s3_q03_1}
\qn{4} Given that mass of Earth is M and that of moon is m and distance between them is R. Neglect effect of any other planet or sun. A small mass is projected from midway from Earth and moon its escape velocity is $P\sqrt{\dfrac GR(m+M)}$ find P.
\qn{5} Snell's law relation for air glass boundary is given by equation $n_{air}\sin\theta_i=n_{glass}\sin\theta_r$. Here angle of refraction in glass is $\theta_r$ and the angle of incident in air is $\theta_i$. The refractive index of air is n$_{air}$= 1 (with negligible uncertainty). In an experiment to verify Snell's Law, the best fit line on a graph between sin$\theta_r$ (on the y-axis) and sin$\theta_i$ (on the x-axis) is given by the equation $y=(6.66\pm0.05)\times10^{-1}x+(0.022\pm0.024)$. The refractive index of the glass, including the error is reported as 1.ab $\pm$ 0.0c, where a, b and c are digits after the decimal point. Calculate value of $(a+b-c)$.
\qn{6} In given circuit diagram initial charge on parallel plates of capacitor C$_1$ is shown in figure. Now all the switches are closed simultaneously. Find the ratio of electrostatic force of attraction between plates of capacitor C$_1$ before closing the switches to steady state reached after closing the switches.\par
\qfig{p2_phy_s3_q06_1}
\partsec{2}{CHEMISTRY}
\sect{SECTION-I(i)}
\qn{1} Potassium ferrocyanide is an important laboratory reagent used to detect many cations. Select the metal ions that can give almost white ppt. with it.
\begin{qoptions}
\opt{A}{Zn$^{2+}$}
\opt{B}{Fe$^{3+}$}
\opt{C}{Cu$^{2+}$}
\opt{D}{Fe$^{2+}$}
\end{qoptions}
\qn{2} Which of the following are ``True'' statements?
\begin{qoptions}
\opt{A}{Zn and Al salts can be distinguished by using NH$_4$OH.}
\opt{B}{Zn and Al salts can be distinguished by using NaOH}
\opt{C}{ZnS is precipitated by H$_2$S in ammonical solution of zinc salt but not in the acidic solution.}
\opt{D}{Al(OH)$_3$ is soluble in NaOH as well as HCl}
\end{qoptions}
\qn{3} The following diketodialdehyde is the result of ozonolysis (O$_3$ followed by Zn) of which diene(s)?\par
\qfig{p2_chem_s1_q03_1}
\begin{qoptions}
\opt{A}{\qfigi{p2_chem_s1_q03_A}}
\opt{B}{\qfigi{p2_chem_s1_q03_B}}
\opt{C}{\qfigi{p2_chem_s1_q03_C}}
\opt{D}{\qfigi{p2_chem_s1_q03_D}}
\end{qoptions}
\qfullbegin
\qn{4} Identify reactions correctly matched with their major products :\par
\qfig{p2_chem_s1_q04_1}
\qfullend
\qn{5} The equilibrium constant (K$_C$) for reaction of weak acid HA with strong base NaOH is 10$^9$ at 25$^\circ$C then -
\begin{qoptions}
\opt{A}{Ionisation constant of HA is 10$^{-5}$ at 25$^\circ$C}
\opt{B}{At 25$^\circ$C, pH of 0.01 M aq. solution of HA is 3.5}
\opt{C}{At 25$^\circ$C, pH of 0.1 M aq. solution of NaA is 9}
\opt{D}{If K$_b$ of weak acid BOH is 10$^{-4}$ at 25$^\circ$C, then equilibrium constant for neutralisation reaction between HA and BOH will be 10$^{-5}$ at 25$^\circ$C.}
\end{qoptions}
\qn{6}\par
\qfig{p2_chem_s1_q06_1}
$\mathrm{E^0_{A/A^+}}=1.6$ volt. \quad $\mathrm{E^0_{B/B^+}}=1.5$ volt.\par
In the given circuit a galvanic cell is connected to external resistance of 2$\Omega$ and external battery of 0.04 volt. If A and B are metal electrodes then $\left(\dfrac{RT}{F}2.303=0.06\right)$
\begin{qoptions}
\opt{A}{Initially current will flow from A to B (outside the solution)}
\opt{B}{Initially not more than 0.03 ampere current will flow in circuit}
\opt{C}{By connecting external battery we are charging galvanic cell}
\opt{D}{When current in circuit is zero, value of $\log\dfrac{[A^+]}{[B^+]}=1$}
\end{qoptions}
\qn{7} During discharging a lead storage battery following reactions occurs.\par
Pb + PbO$_2$ + 2H$_2$SO$_4$ $\to$ 2PbSO$_4$ + 2H$_2$O\par
[Atomic mass : Pb = 207, S = 32]
\begin{qoptions}
\opt{A}{For production of 606 milligram of PbSO$_4$, 4 milli equivalents of H$_2$SO$_4$ are required}
\opt{B}{On producing 606 milligram of PbSO$_4$ , 2 milli equivalents of H$_2$O are produced}
\opt{C}{If 4 millifaraday electricity is drawn from cell at potential of 2 volt then work done by cell is 772 Joules ideally}
\opt{D}{During discharging density of H$_2$SO$_4$ solution decreases.}
\end{qoptions}
\qn{8} During electrolysis of aq.CuBr$_2$ using Pt electrode -
\begin{qoptions}
\opt{A}{Br$_2$ produced at anode}
\opt{B}{Cu produced at cathode}
\opt{C}{H$_2$ produced at cathode}
\opt{D}{O$_2$ produced at anode}
\end{qoptions}
\sect{SECTION-I(ii)}
\qn{9} \textbf{Answer the following by appropriately matching the lists based on the information given in the paragraph.}\par
\smallskip
\noindent{\footnotesize\setlength{\tabcolsep}{2pt}\begin{tabular}{|L{.8cm}|L{2.0cm}|L{.8cm}|L{3.3cm}|}\hline
\multicolumn{2}{|c|}{List-I: Radicals} & \multicolumn{2}{c|}{\parbox{3.9cm}{\centering List-II: Correct statement regarding them}}\\\hline
(I) & Br$^-$ & (P) & Can be confirmed by Brown ring test\\\hline
(II) & I$^-$ & (Q) & Can be confirmed by Layer test\\\hline
(III) & S$^{2-}$ & (R) & Can be confirmed by methylene blue test\\\hline
(IV) & NO$_3^-$ & (S) & On reaction with conc. H$_2$SO$_4$, oxidation takes place\\\hline
 & & (T) & All salts are soluble\\\hline
 & & (U) & On reaction with dilute H$_2$SO$_4$ produces colourless gas\\\hline
\end{tabular}}\par
Which of the following option has the \textbf{CORRECT} combination considering List-I \& List-II ?
\begin{qoptions}
\opt{A}{I, P}
\opt{B}{II, T}
\opt{C}{III, R}
\opt{D}{IV, S}
\end{qoptions}
\qn{10} \textbf{Answer the following by appropriately matching the lists based on the information given in the paragraph.}\par
\smallskip
\noindent{\footnotesize\setlength{\tabcolsep}{2pt}\begin{tabular}{|L{.8cm}|L{2.0cm}|L{.8cm}|L{3.3cm}|}\hline
\multicolumn{2}{|c|}{List-I: Radicals} & \multicolumn{2}{c|}{\parbox{3.9cm}{\centering List-II: Correct statement regarding them}}\\\hline
(I) & Br$^-$ & (P) & Can be confirmed by Brown ring test\\\hline
(II) & I$^-$ & (Q) & Can be confirmed by Layer test\\\hline
(III) & S$^{2-}$ & (R) & Can be confirmed by methylene blue test\\\hline
(IV) & NO$_3^-$ & (S) & On reaction with conc. H$_2$SO$_4$, oxidation takes place\\\hline
 & & (T) & All salts are soluble\\\hline
 & & (U) & On reaction with dilute H$_2$SO$_4$ produces colourless gas\\\hline
\end{tabular}}\par
Which of the following is the \textbf{INCORRECT} combination for List-I and List-II.
\begin{qoptions}
\opt{A}{II, S}
\opt{B}{IV, Q}
\opt{C}{I, Q}
\opt{D}{III, U}
\end{qoptions}
\qn{11} \textbf{Answer the following by appropriately matching the lists based on the information given in the paragraph.}\par
\smallskip
\noindent{\footnotesize\setlength{\tabcolsep}{2pt}\begin{tabular}{|L{.8cm}|L{2.0cm}|L{.8cm}|L{3.3cm}|}\hline
\multicolumn{2}{|c|}{List-I} & \multicolumn{2}{c|}{List-II}\\\hline
(I) & Acetophenone & (P) & Reaction with AgNO$_3$ + NH$_4$OH to give white ppt.\\\hline
(II) & Benzaldehyde & (Q) & Reaction with NaOI to give yellow ppt. of CHI$_3$\\\hline
(III) & Butanal & (R) & Reaction with Fehling to give brown ppt. of Cu$_2$O\\\hline
(IV) & 1-butyne & (S) & Reaction with AC$_2$O + AcONa to give Cinnamic acid\\\hline
 & & (T) & Reaction with H$_2$N-OH/H$^+$ gives 2-oxime products\\\hline
 & & (U) & Reaction with Tollen's reagent give silver mirror\\\hline
\end{tabular}}\par
The correct code for your answer :-
\begin{qoptions}
\opt{A}{I $\to$ U; II $\to$ P,Q,R; III $\to$ R,S,T; IV $\to$ P,S,U}
\opt{B}{I $\to$ R; II $\to$ R,S; III $\to$ P,Q; IV $\to$ T,U}
\opt{C}{I $\to$ T; II $\to$ T,U; III $\to$ T,U; IV $\to$ P}
\opt{D}{I $\to$ R,S,T; II $\to$ P,Q,R; III $\to$ P,Q; IV $\to$ R,S,T}
\end{qoptions}
\qn{12} \textbf{Answer the following by appropriately matching the lists based on the information given in the paragraph.}\par
\smallskip
\noindent{\footnotesize\setlength{\tabcolsep}{2pt}\begin{tabular}{|L{.8cm}|L{2.0cm}|L{.8cm}|L{3.3cm}|}\hline
\multicolumn{2}{|c|}{List-I} & \multicolumn{2}{c|}{List-II}\\\hline
(I) & Acetophenone & (P) & Reaction with AgNO$_3$ + NH$_4$OH to give white ppt.\\\hline
(II) & Benzaldehyde & (Q) & Reaction with NaOI to give yellow ppt. of CHI$_3$\\\hline
(III) & Butanal & (R) & Reaction with Fehling to give brown ppt. of Cu$_2$O\\\hline
(IV) & 1-butyne & (S) & Reaction with AC$_2$O + AcONa to give Cinnamic acid\\\hline
 & & (T) & Reaction with H$_2$N-OH/H$^+$ gives 2-oxime products\\\hline
 & & (U) & Reaction with Tollen's reagent give silver mirror\\\hline
\end{tabular}}\par
The correct code for your answer :-
\begin{qoptions}
\opt{A}{I $\to$ Q; II $\to$ S; III $\to$ R; IV $\to$ P}
\opt{B}{I $\to$ P; II $\to$ S; III $\to$ R; IV $\to$ Q}
\opt{C}{I $\to$ Q; II $\to$ R; III $\to$ S; IV $\to$ P}
\opt{D}{I $\to$ P,S,T; II $\to$ Q; III $\to$ S; IV $\to$ R}
\end{qoptions}
\sect{SECTION-II}
\qn{1} Identify total number of carboxylic acids with C$_5$H$_{10}$O$_2$ molecular formula gives 1$^\circ$-Alcohol on reaction with LiAlH$_4$.
\qn{2} Compare the given characteristics of aniline and cyclohexanamine :\par
(i) Both are primary amine\par
(ii) Both can be acylated by RCOCl\par
(iii) Both reacts with CHCl$_3$/KOH\par
(iv) Both reacts with NaNO$_2$ + HCl at 0-5$^\circ$C\par
(v) Both reacts with PhSO$_2$Cl to give a compound which is soluble in KOH\par
(vi) Both gives coupling reaction with phenol\par
(vii) Both gives electrophilic substitution reaction\par
How many of the given characteristics are correct?
\qn{3} Total number of anions which will give coloured gas/vapour will be :-\par
\qfig{p2_chem_s3_q03_1}
$\mathrm{NO_2^{\ominus}},\mathrm{S_2O_3^{2-}},\mathrm{Cl^-},\mathrm{Br^-},\mathrm{I^-},\mathrm{NO_3^-},\mathrm{F^-}$
\qn{4} How many of the following is/are yellow and water insoluble substances ?\par
Ag$_2$CrO$_4$, \quad PbCrO$_4$, \quad Na$_2$CrO$_4$, \quad BaCrO$_4$, \quad CrO$_2$Cl$_2$, \quad (NH$_4$)$_2$Cr$_2$O$_7$, \quad (NH$_4$)$_3$PO$_4$.12MoO$_3$, \quad K$_3$[Co(NO$_2$)$_6$], \quad AgI, \quad PbI$_2$
\qn{5} A saturated solution of B(OH)$_3$ is prepared in pure water at 25$^\circ$C if Ksp of B(OH)$_3$ is 2.7 $\times$ 10$^{-33}$ then concentration of B$^{+3}$ is y $\times$ 10$^{-12}$ calculate value of y will be -\par
B(OH)$_3$(s) $\rightleftharpoons$ B$^{+3}$(aq.) + 3OH$^-$(aq.)
\qn{6} KMnO$_4$ + HCl $\to$ Cl$_2$ + KCl + MnCl$_2$ + H$_2$O\par
Due to reaction if 162 gm H$_2$O are produced then how many moles of MnCl$_2$ will be obtained?
\partsec{3}{MATHEMATICS}
\sect{SECTION-I(i)}
\qn{1} If $20\sin^2B-6\sin A+8\cos A-20\sin B+15=0$, then -
\begin{qoptions}
\opt{A}{$\angle A=(2n+1)\dfrac\pi2+\tan^{-1}\dfrac34$, (n $\in$ I)}
\opt{B}{$\angle A=(4n+1)\dfrac\pi2+\tan^{-1}\dfrac43$, (n $\in$ I)}
\opt{C}{$\angle B=\dfrac\pi6$}
\opt{D}{$\angle B=\dfrac\pi3$}
\end{qoptions}
\qn{2} The number of values of $\alpha$ in [0, $3\pi$] satisfying the equation $|\sin4x|+|\cos4x|=|\cos2\alpha|$ is A and sum of all values of $\alpha$ is B then which of the following option(s) is/are correct-
\begin{qoptions}
\opt{A}{A + 2B = 7 + 21$\pi$}
\opt{B}{A + 2B = 6 + 21$\pi$}
\opt{C}{2AB = 147$\pi$}
\opt{D}{2AB = 21$\pi$}
\end{qoptions}
\qn{3} If $16+\sin^{12}x+\cos^{11}x-\sec^2y-9\operatorname{cosec}^2y\ \ge\ \sec^8z-2\sec z+\log_2\left(5-\cos^2x\right)$ and x,y,z $\in$ [0,2$\pi$] then correct option/s is/are-
\begin{qoptions}
\opt{A}{Number of possible values of x are 2}
\opt{B}{Number of possible ordered pairs(x,y) are 8}
\opt{C}{Number of possible ordered pairs (y,z) are 8}
\opt{D}{Number of possible ordered triplets (x,y,z) are 16}
\end{qoptions}
\qn{4} If $\cos A=3\cos(A+2B)$, then the value of $\cot(A+B)\cot B$ is -
\begin{qoptions}
\opt{A}{odd number}
\opt{B}{even number}
\opt{C}{prime number}
\opt{D}{composite number}
\end{qoptions}
\qn{5} Let the radius of circumcircle of an isosceles triangle ABC is equal to AB = AC. If $P=\sin^2A+\sin^2B$ and $Q=\dfrac rR$, then (where r and R are inradius and circumradius of $\Delta$ABC respectively)-
\begin{qoptions}
\opt{A}{P = 1}
\opt{B}{$P=\dfrac12$}
\opt{C}{$Q=\dfrac{2\sqrt3-3}{2}$}
\opt{D}{$Q=\dfrac{3\sqrt3-2}{2}$}
\end{qoptions}
\qn{6} If in $\Delta$ABC, sides a,b,c are in AP and $\cos C=\dfrac13$, then (notations have their usual meanings)-
\begin{qoptions}
\opt{A}{$\Delta$ABC is obtuse angle triangle}
\opt{B}{$\tan\left(\dfrac{B-A}{2}\right)=\dfrac{\sqrt2}{19}$}
\opt{C}{2Rr = 33 for a = 9}
\opt{D}{$\sin\dfrac A2\sin\dfrac B2\sin\dfrac C2=\dfrac{4}{33}$}
\end{qoptions}
\qn{7} In $\Delta$ABC, $3\sin^2B+4\sin B\sin A+16\sin^2A-8\sin B-20\sin A+9=0$, then -
\begin{qoptions}
\opt{A}{$\sin^2A+\sin^2B+\sin^2C=2$}
\opt{B}{$\cos2A+\cos2B+\cos2C=-1$}
\opt{C}{$\cot A+\cot B+\cot C=\sqrt3$}
\opt{D}{$\sin2A+\sin2B+\sin2C=\sqrt3$}
\end{qoptions}
\qn{8} Let $f(n)=\displaystyle\sum_{r=1}^{n}\dfrac{\cos\left(2\times3^{r-1}\alpha\right)}{\sin\left(3^r\alpha\right)}$, where $\alpha=18^\circ$ and $g(n)=\displaystyle\sum_{r=1}^{n}\dfrac{\cos\left(3\times2^{r-1}\beta\right)}{\cos\left(2^r\beta\right)}$ where $\beta=36^\circ$, then
\begin{qoptions}
\opt{A}{$f(2020)+g(2020)=2018$}
\opt{B}{$f(2020)+g(2020)=2020$}
\opt{C}{$f(2021)-g(2021)=-2018$}
\opt{D}{$f(2021)-g(2021)=-2019$}
\end{qoptions}
\sect{SECTION-I(ii)}
\qn{9} \textbf{Answer the following by appropriately matching the lists based on the information given in the paragraph.}\par
Let $a=\cot^2\dfrac\pi7+\cot^2\dfrac{2\pi}7+\cot^2\dfrac{3\pi}7$, b = sum of all possible values of x when\par
$\operatorname{cosec}x^\circ=\dfrac{1}{\sin45^\circ.\sin46^\circ}+\dfrac{1}{\sin47^\circ.\sin48^\circ}+\ldots+\dfrac{1}{\sin133^\circ\sin134^\circ}$ where $x^\circ\in[0^\circ,360^\circ]$\par
$c=\dfrac{1}{\sqrt2\sin5^\circ\sin15^\circ\sin25^\circ\ldots.\sin85^\circ}$ and $d=\tan^210^\circ+\tan^250^\circ+\tan^270^\circ$.\par
There are some expression given in the list-I whose values are given in list-II below.\par
\smallskip
\noindent{\footnotesize\setlength{\tabcolsep}{2pt}\begin{tabular}{|L{.95cm}|L{3.2cm}|L{.95cm}|L{1.8cm}|}\hline
\multicolumn{2}{|c|}{List-I} & \multicolumn{2}{c|}{List-II}\\\hline
(I) & $\dfrac{b+c+d}{a}$ & (P) & 18\\\hline
(II) & $\dfrac{c+d-b}{a}$ & (Q) & 17\\\hline
(III) & $\dfrac bd$ & (R) & 64\\\hline
(IV) & $\dfrac{c}{d-a}$ & (S) & 60\\\hline
 & & (T) & 89\\\hline
 & & (U) & 20\\\hline
\end{tabular}}\par
Which of the following is the only \textbf{CORRECT} combination ?
\begin{qoptions}
\opt{A}{(IV), (S)}
\opt{B}{(III), (Q)}
\opt{C}{(II), (U)}
\opt{D}{(I), (T)}
\end{qoptions}
\qn{10} \textbf{Answer the following by appropriately matching the lists based on the information given in the paragraph.}\par
Let $a=\cot^2\dfrac\pi7+\cot^2\dfrac{2\pi}7+\cot^2\dfrac{3\pi}7$, b = sum of all possible values of x when\par
$\operatorname{cosec}x^\circ=\dfrac{1}{\sin45^\circ.\sin46^\circ}+\dfrac{1}{\sin47^\circ.\sin48^\circ}+\ldots+\dfrac{1}{\sin133^\circ\sin134^\circ}$ where $x^\circ\in[0^\circ,360^\circ]$\par
$c=\dfrac{1}{\sqrt2\sin5^\circ\sin15^\circ\sin25^\circ\ldots.\sin85^\circ}$ and $d=\tan^210^\circ+\tan^250^\circ+\tan^270^\circ$.\par
There are some expression given in the list-I whose values are given in list-II below.\par
\smallskip
\noindent{\footnotesize\setlength{\tabcolsep}{2pt}\begin{tabular}{|L{.95cm}|L{3.2cm}|L{.95cm}|L{1.8cm}|}\hline
\multicolumn{2}{|c|}{List-I} & \multicolumn{2}{c|}{List-II}\\\hline
(I) & $\dfrac{b+c+d}{a}$ & (P) & 18\\\hline
(II) & $\dfrac{c+d-b}{a}$ & (Q) & 17\\\hline
(III) & $\dfrac bd$ & (R) & 64\\\hline
(IV) & $\dfrac{c}{d-a}$ & (S) & 60\\\hline
 & & (T) & 89\\\hline
 & & (U) & 20\\\hline
\end{tabular}}\par
Which of the following is the only \textbf{CORRECT} combination ?
\begin{qoptions}
\opt{A}{(IV), (U)}
\opt{B}{(IV), (P)}
\opt{C}{(II), (Q)}
\opt{D}{(II), (S)}
\end{qoptions}
\qn{11} \textbf{Answer the following by appropriately matching the lists based on the information given in the paragraph.}\par
If set $X=\left\{x:\dfrac{\sqrt3-1}{\sin x}+\dfrac{\sqrt3+1}{\cos x}=4\sqrt2,\right.$\par $\left.x\in(0,\pi)-\left\{\dfrac\pi2\right\}\right\}$\par
$Y=\left\{x:4|\sin x|-|x-2|=0\right\}$\par
$Z=\left\{f(x):f(x)=\sin(\sin x)+\cos(\sin x)\right\}$\par
$W=\left\{f(x):f(x)=4\sin x\cos x-\sin^2x+3\cos^2x\right\}$\par
\smallskip
\noindent{\footnotesize\setlength{\tabcolsep}{2pt}\begin{tabular}{|L{.95cm}|L{3.2cm}|L{.95cm}|L{1.8cm}|}\hline
\multicolumn{2}{|c|}{List-I} & \multicolumn{2}{c|}{List-II}\\\hline
(I) & Number of elements in set X & (P) & 1\\\hline
(II) & Number of elements in set Y & (Q) & 2\\\hline
(III) & Number of integers in set Z & (R) & 3\\\hline
(IV) & Number of integers in set W & (S) & 4\\\hline
 & & (T) & 5\\\hline
 & & (U) & 6\\\hline
\end{tabular}}\par
Which of the following is the only \textbf{INCORRECT} combination ?
\begin{qoptions}
\opt{A}{(II), (U)}
\opt{B}{(III), (P)}
\opt{C}{(I), (R)}
\opt{D}{(IV), (T)}
\end{qoptions}
\qn{12} \textbf{Answer the following by appropriately matching the lists based on the information given in the paragraph.}\par
If set $X=\left\{x:\dfrac{\sqrt3-1}{\sin x}+\dfrac{\sqrt3+1}{\cos x}=4\sqrt2,\right.$\par $\left.x\in(0,\pi)-\left\{\dfrac\pi2\right\}\right\}$\par
$Y=\left\{x:4|\sin x|-|x-2|=0\right\}$\par
$Z=\left\{f(x):f(x)=\sin(\sin x)+\cos(\sin x)\right\}$\par
$W=\left\{f(x):f(x)=4\sin x\cos x-\sin^2x+3\cos^2x\right\}$\par
\smallskip
\noindent{\footnotesize\setlength{\tabcolsep}{2pt}\begin{tabular}{|L{.95cm}|L{3.2cm}|L{.95cm}|L{1.8cm}|}\hline
\multicolumn{2}{|c|}{List-I} & \multicolumn{2}{c|}{List-II}\\\hline
(I) & Number of elements in set X & (P) & 1\\\hline
(II) & Number of elements in set Y & (Q) & 2\\\hline
(III) & Number of integers in set Z & (R) & 3\\\hline
(IV) & Number of integers in set W & (S) & 4\\\hline
 & & (T) & 5\\\hline
 & & (U) & 6\\\hline
\end{tabular}}\par
Which of the following is the only \textbf{CORRECT} combination ?
\begin{qoptions}
\opt{A}{(IV), (U)}
\opt{B}{(IV), (T)}
\opt{C}{(I), (P)}
\opt{D}{(I), (S)}
\end{qoptions}
\sect{SECTION-II}
\qn{1} Let $\displaystyle\prod_{r=1}^{51}\tan\left(\dfrac\pi3\left(1+\dfrac{3^r}{3^{50}-1}\right)\right)=k\prod_{r=1}^{51}\cot\left[\dfrac\pi3\left(1-\dfrac{3^r}{3^{50}-1}\right)\right]$.\par
On solving equation we get, $1-3\tan^2\left(\dfrac\pi{3^{50}-1}\right)=\dfrac{a}{bk-1}$, (a,b $\in$ I) then value of (a $-$ b) is equal
\qn{2} With usual notations, if in a $\Delta$ABC, a = 6 \& bcotB + ccotC is equal to twice of sum of its inradius \& its circumradius, then its maximum area is
\qn{3} In $\Delta$ABC, if $\dfrac{\sum(b+c)\cos A}{4\cos\frac A2\cos\frac B2\cos\frac C2}=8$, then the value of circumradius of $\Delta$ABC, is
\qn{4} Find the numerical value of $\operatorname{cosec}^2\dfrac\pi7+\operatorname{cosec}^2\dfrac{2\pi}7+\operatorname{cosec}^2\dfrac{4\pi}7$.
\qn{5} In acute angled triangle ABC with usual notations, perpendicular from circumcentre upon side BC intersects circumcircle at D, AD = $\sqrt{10}$, $\sin(B-C)=\dfrac35$, then $\dfrac1R$ is (where R is circumradius of triangle ABC)
\qn{6} If $\tan\alpha=\dfrac{x^2-x}{x^2-x+1}$, $\tan\beta=\dfrac{1}{2x^2-2x+1}$, $0<\alpha,\beta<\dfrac\pi2$ so that $\tan\left(\dfrac{\alpha+\beta}3\right)=a-\sqrt b$, a and b are coprime, then a + b is equal to
\end{multicols}
\end{document}
